\section{Read the function interfaces independently}\label{ball-x-interfaces}
\Needspace{17\baselineskip}
\tagpdfsetup{table/header-rows={1}}
\begin{longtable}{p{32mm}p{59mm}p{56mm}}
  \toprule
  Function                   & Arguments                                                                            & Results and permitted effects                                                                            \\
  \midrule
  \texttt{advance\_axis}     & \texttt{\$a0} raw position; \texttt{\$a1} raw velocity                               & \texttt{\$v0} position; \texttt{\$v1} velocity; provided static objects unchanged                        \\
  \texttt{render\_ball}      & \texttt{\$a0} destination; \texttt{\$a1}, \texttt{\$a2} raw x, y; \texttt{\$a3} mode & \texttt{\$v0} pixel ID; destination frame written; other provided objects unchanged                      \\
  \texttt{add\_fixed\_probe} & \texttt{\$a0}, \texttt{\$a1} signed word inputs as full register values              & \texttt{\$v0} sign-extended wrapped word; \texttt{\$v1} exact doubleword sum; provided objects unchanged \\
  \bottomrule
\end{longtable}
All three preserve the full contents of \texttt{\$s0} through \texttt{\$s7}
and restore \texttt{\$sp}. The caller can supply different return continuations,
so return through the link value rather than a global exit label. Temporary and
argument registers may change. The probe's operand domain is wider than the
motion domain; keep those numeric promises separate.
